\section{Введение: ролевые обозначения и границы вызовов}
\label{sec:dev40-context}

В конвейерах генерации кода обычно выделяют этапы планирования, реализации и проверки. В таком описании часто скрыты два различных воздействия. Ролевое обозначение может изменить инструкции, данные одному ответу модели, тогда как граница вызова может разделить работу между отдельно сэмплированными ответами с передачей промежуточного текста. Если оба воздействия изменяются одновременно, наблюдаемое различие нельзя приписать ни формулировке ролей, ни повторному выводу, ни информации, раскрываемой между вызовами. Настоящее исследование на этапе разработки изолирует этот вопрос на уровне протокола: отличаются ли инструкции с ролевыми обозначениями от нейтральных инструкций этапов при изменении топологии вызовов в одной и той же совокупности задач и при одном и том же предоставленном контексте?

Оцениваемый эффект намеренно уже, чем утверждение о независимых агентах или скрытом мышлении. В исследовании сравниваются пять реализованных условий: прямой решатель с одним вызовом, ответ с нейтральными этапами в одном вызове, ответ с ролевыми этапами в одном вызове, нейтральный конвейер из трёх вызовов и конвейер из трёх вызовов с ролевыми обозначениями. В условиях одного вызова три процедурных этапа размещаются в одном ответе модели. В условиях нескольких вызовов каждый следующий этап получает предыдущий ответ целиком. Поэтому топология вызовов и поток промежуточной информации являются частью определения воздействия. Результат может поддержать утверждение об этих зафиксированных промптах и границах, но не может установить, что профессиональные обозначения создают независимых внутренних агентов.

\section{Связанные работы}
\label{sec:dev40-related}

Ролевое описание имеет предшественников в программном промптинге и внутреннем диалоге. INoT организует рассуждение с помощью явных этапов и инструкций~\cite{sun2025inot}. Эта работа мотивирует контролируемое сравнение, однако само наличие названий ролей не определяет причинный эффект ролей. Полезное сравнение должно отделять формулировку ролей от числа вызовов модели и сохранять информацию, доступную каждому условию.

Бюджет и структура взаимодействия также представляют собой разные факторы в связанных исследованиях. Работа об односоставных и многосоставных рассуждающих системах при сопоставимых бюджетах токенов мышления подчёркивает, что дополнительные вызовы и дополнительные вычисления не являются взаимозаменяемыми воздействиями~\cite{tran2026budget}. Поэтому в настоящем исследовании фиксируются счётчики входных, кэшированных входных и выходных токенов, а оценка по стандартным API-тарифам приводится как описательная мера ресурса. Канал подписки не трактуется как счёт API, а сравнение токенов не объявляется полным сравнением совокупной стоимости владения.

В области программной инженерии Agentless использует структурированный рабочий процесс локализации, исправления и проверки~\cite{xia2024agentless}, тогда как SWE-agent использует интерфейс репозитория с инструментами~\cite{yang2024sweagent}. Эти системы предназначены для практического исправления репозиториев и включают политики поиска или взаимодействия, отсутствующие в используемых здесь задачах функций с фиксированным контекстом. BigCodeBench предоставляет разнообразные задачи генерации кода с вызовами библиотек и сложными инструкциями~\cite{zhuo2024bigcodebench}. Он подходит для проверки следования инструкциям и исполнимых результатов, но распределение задач на уровне функций не заменяет исправление репозиториев. SWE-bench оценивает изменения по тестам репозитория~\cite{jimenez2023swebench,swebenchdocs2026} и остаётся отдельной внешней границей настоящего исследования.

\section{Выполненные методы исследования на этапе разработки}
\label{sec:dev40-methods}

Выполненное расширение использует GPT-5.4 mini со средним уровнем рассуждения через Codex CLI 0.153.4. Решатель работает через канал подписки и зафиксированный псевдоним модели; утверждение о неизменном снимке весов не делается. Каждое условие получает один и тот же текст задачи и полный предоставленный контекст. В условиях нескольких вызовов каждый последующий этап получает весь предшествующий раскрытый ответ. Вход не сжимается и не сокращается выборочно для отдельного условия. Во время генерации отключены инструменты, просмотр веб-страниц, делегирование и выполнение тестов; для каждого свежего эфемерного потока требуется строгая трасса JSONL.

Совокупность задач состоит из 40 случайно упорядоченных задач разработки из BigCodeBench v0.1.4. Порядок, заданный начальным числом, и точное распределение задач записаны в неизменяемом манифесте. Все задачи сохраняются, включая задачу, контроль эталона которой непригоден для измерения корректности. Две помеченные повторами репликации, с идентификаторами 2 и 3, сопоставляются внутри задачи и условия. Эти обозначения являются идентификаторами повторов, а не сэмплирующими сидами провайдера: CLI не предоставляет проверяемый сид сэмплирования, а температура не контролируется. Следовательно, повторы дают повторные наблюдения в рамках одного заявленного протокола, но не означают детерминированного воспроизведения модели.

Пять условий, применённые к 40 задачам и двум меткам повторов, дают 400 назначенных программ-кандидатов. В условиях одного вызова и прямого решения на одного кандидата приходится один CLI-вызов; каждое трёхэтапное условие использует три свежих вызова. Поэтому полная матрица содержит 720 свежих CLI-вызовов. Упорядоченные задачи разделяются на четыре непересекающиеся группы по десять задач для каждой волны повторов. Четыре группы запускаются параллельно, каждая с двумя рабочими потоками решателя, что даёт не более восьми одновременно работающих CLI-процессов. Волна повторов 3 начинается только после завершения волны 2. Порядок отправки внутри каждой группы рандомизируется зафиксированным решателем, тогда как состав задач и обозначения условий остаются заданными манифестом.

Решатель устанавливает ограничение в 180 секунд на один CLI-вызов, предварительную проверку входа на 65\,536 байт и отсутствие единого жёсткого лимита завершения, сопоставленного между условиями. Отсутствие сопоставленного лимита является заявленным ограничением: дизайн фиксирует информацию о задаче и настройки, не относящиеся к воздействию, но не фактическую длину завершения, использование рассуждения, время выполнения или общие расходы. Автоматические повторы, исправление архивированного итогового кандидата после выполнения, сжатие раскрытых ответов перед передачей и выбор более удачного блока кода запрещены. Этапы реализации и проверки внутри заявленного воздействия могут пересматривать код согласно инструкции. Если группа завершается ошибкой, отправленные события и наблюдаемое использование сохраняются, уже активная волна доводится до конца, а следующая волна повторов не запускается. Частичная матрица остаётся частичной, а ненаблюдавшиеся назначенные ячейки сохраняются как пропуски.

До генерации модели эталонные и намеренно неправильные контроли были выполнены для всех 40 назначенных задач через неизменённый upstream CLI BigCodeBench в итоговой изолированной среде, с исходными тестами, отключённой сетью и штатным режимом ``instruct''. Контрольные данные показывают 39 задач с успешно прошедшим эталоном и не прошедшим неправильным решением. Для одной задачи эталон завершился ошибкой, поэтому она сохраняется в назначении, но исключается из измерения корректности во всех условиях. Архивированы gate, хэши образа и пакетов, набора данных и подготовленного разделения, дерева исходников, команд и исходных контрольных отчётов. Прохождение контроля подтверждает исполнимость и различение неправильного решения, но не удостоверяет семантическую корректность каждого утверждения бенчмарка.

Независимой единицей является задача. Повторы усредняются внутри задачи только для парных описательных сравнений, когда наблюдаемы оба повтора и оба сравниваемых условия. Анализ сообщает доли среди оцениваемых попыток вместе с границами пропусков по полному назначенному знаменателю. Сравнения ресурсов используют полные парные средние по задачам; описательные средние ресурсов по отдельным условиям сохраняют знаменатели всех наблюдаемых кандидатов, когда матрица неполна. Бутстрэп-интервалы по кластерам задач строятся по 10\,000 ресэмплированиям с фиксированным начальным числом анализа. Это описательные 95\% бутстрэп-интервалы без подтверждающего статистического теста или расчётного утверждения о не меньшей эффективности; 40 задач не превращаются в 400 независимых наблюдений только потому, что матрица содержит 400 кандидатов.

\section{Ограничения и внешняя граница}
\label{sec:dev40-limits}

Выполненный дизайн характеризует поведение одного запрошенного псевдонима модели, одной настройки рассуждения, одной версии CLI, пяти условий промптинга и выбранной совокупности из 40 задач разработки. Он не характеризует устойчивое свойство модели у провайдера, универсальный эффект ролей или причинный эффект, независимый от потока информации. Поскольку сид сэмплирования провайдера и контроль температуры недоступны, метки повторов следует интерпретировать как повторные наблюдения на уровне протокола, а не как точные стохастические репликации. Отсутствие общего жёсткого лимита вывода оставляет качество, длину завершения и использование ресурсов совместно информативными, но не уравненными по бюджету.

Валидность тестов BigCodeBench остаётся отдельным вопросом. Независимый аудит инструкций и тестов фиксирует проблемы согласованности промптов и тестов как аннотации, тогда как исходные промпты, тесты и оценки сохраняются без изменений. Успешный эталонный контроль необходим для включения задачи в оценивание, но не может доказать, что каждое утверждение измеряет естественно-языковую задачу. Все назначенные задачи остаются в знаменателе. Контрольные или инфраструктурные ошибки дают пропущенные результаты; расхождения инструкций и тестов остаются аннотациями и не стирают исходные оценки. Ресэмплирование задач учитывает повторные наблюдения внутри задачи, но не устраняет зависимость между связанными задачами бенчмарка.

Исследование также ограничено относительно утверждений о внешних агентах для репозиториев. В нём не выполняются нативные системы INoT, Agentless или SWE-agent и не приводится внешний результат этих систем на бенчмарке. Эти системы различаются поиском, использованием инструментов, состоянием репозитория, применением патчей и взаимодействием с тестами. Будущее сравнение должно зафиксировать каждую нативную реализацию, модель и бюджет, собрать полные траектории и использовать соответствующий официальный оценщик. В частности, результат генерации функции нельзя переименовать в свидетельство о ремонте репозитория, а опубликованный внешний процент нельзя вставить в эту матрицу как парное условие.

Таким образом, настоящее исследование представляет собой ограниченный выполненный эксперимент на этапе разработки. Оно даёт проверяемое сравнение ролевых обозначений и границ вызовов при полном предоставленном контексте, повторно раскрываемой истории, нативных тестах BigCodeBench и явном учёте ресурсов. Само по себе оно не обосновывает утверждение о скрытом многоагентном поведении, превосходстве на всём бенчмарке, не меньшей эффективности или внешней валидности для современной разработки репозиториев.
